\documentclass[exercice,a4]{thedocument}%
\usepackage[french]{babel}%
\usepackage{enumitem}%
\usepackage[np]{numprint}%
\startdatakeys
    \source{21GENMATNC1}
    \BO{2026}
    \cycle{4}
    \competencesDuSocle{CA,RA}
    \objectifsApprentissage{C4A1aa,C4A3aa}
    \motsCles{facteurs premiers,nombre premier,pourcentage,PGCD,multiple}
\enddatakeys
\begin{document}%

\medskip

\begin{enumerate}
\item
	\begin{enumerate}
		\item Justifier que $330$ n'est pas un nombre premier.

La décomposition en produit de facteurs premiers de $500$ est: $500 = 2^2 \times 5^3$.
		\item Décomposer $330$ en produit de facteurs premiers.
		\item Justifier que $165$ divise $330$.
		\item Justifier que 165 ne divise pas $500$.
	\end{enumerate}
\end{enumerate}

La pâtisserie Délices a préparé $330$ biscuits aux noix et $500$ biscuits au chocolat.

\textbf{La pâtisserie souhaite répartir le plus de biscuits possible dans $165$ boites.}

La pâtisserie met le même nombre de biscuits aux noix dans chaque boîte.

\begin{enumerate}[resume]
\item Combien de biscuits aux noix y a-t-il dans chaque boîte ?

La pâtisserie met aussi le même nombre de biscuits au chocolat dans chaque boîte.
\item
	\begin{enumerate}
		\item Combien de biscuits au chocolat y a-t-il dans chaque boîte ?
		\item Combien de biscuits au chocolat reste-t-il ?
	\end{enumerate}

Une boîte de biscuits coûte \np{3650} francs.

À partir de 10 boîtes achetées, la pâtisserie Délices offre une réduction de 5\,\% sur le montant total.
\item Combien va-t-on payer pour l'achat de 12 boîtes ?

Faire apparaÎtre les calculs effectués.

\end{enumerate}

\bigskip

\showthesource
\end{document}
